Jedným zo základných problémov posilňovaného učenia je credit assignment, teda určenie príspevku jednotlivých akcií k výslednej odmene. Tento problém je obzvlášť výrazný v úlohách s terminálnou odmenou, kde je celý epizód hodnotený jediným skalárnym signálom a priebežná spätná väzba chýba. V takýchto podmienkach algoritmus nedostáva informáciu o tom, ktoré akcie prispeli k výsledku pozitívne a ktoré negatívne, čo vedie k vysokej variancii gradientov a výrazne sťažuje učenie, najmä pri dlhých sekvenciách akcií.

Nedávno navrhnutá metóda SCAR (Shapley Credit Assignment for More Efficient RLHF) formuluje generovanie sekvencií ako kooperatívnu hru a využíva Shapleyho hodnoty na rozklad terminálnej odmeny na príspevky jednotlivých tokenov. V kontexte RLHF to umožňuje nahradiť jednu riedku odmenu hustou sekvenciou tokenových odmien, ktoré môžu byť použité na stabilnejšie a efektívnejšie doúčenie veľkých jazykových modelov.

V tejto práci sa myšlienka Shapley-based credit assignment prenáša z RLHF do algoritmu Dreamer. Na rozdiel od RLHF, kde hodnotenie odmeny vyžaduje samostatný reward model natrénovaný na ľudských preferenciách, Dreamer už obsahuje učený model dynamiky prostredia a odmeny. To umožňuje vyhodnocovať kontrafaktuálne varianty trajektórií pomocou latentných imaginárnych rolloutov bez dodatočnej interakcie s reálnym prostredím.

Práca sa zameriava na textové a textovo-procedurálne úlohy, v ktorých agent vykonáva sekvenciu diskrétnych akcií (napríklad generovanie textu, operácie s pamäťou alebo volanie modulov) a odmenu dostáva až na konci epizódy. Po ukončení epizódy je terminálna odmena rozdelená medzi jednotlivé akcie pomocou Shapleyho samplovania. To sa realizuje odhadom očakávanej terminálnej odmeny pre viaceré modifikované varianty tej istej trajektórie, získané nahradením podmnožiny akcií bazálnou akciou, a ich vyhodnotením prostredníctvom latentných rolloutov modelu prostredia Dreamer. Výsledná sekvencia príspevkov akcií je následne použitá ako hustý signál odmeny na tréning politiky aj modelu dynamiky.

V experimentálnej časti je algoritmus Dreamer porovnaný v troch tréningových režimoch: s použitím iba terminálnej odmeny, s ručným reward shapingom a s rozkladom odmeny pomocou Shapleyho hodnôt. Hodnotenie sa zameriava na rýchlosť učenia, stabilitu tréningu a kvalitu naučenej politiky, ako aj na vplyv počtu Shapleyho vzoriek a stratégie segmentácie akcií.